\section{Introduction}
\IEEEPARstart{L}{iDAR} object detection supplies metric scene information for autonomous vehicles and mobile robots. The detector must recognize objects with different footprints and orientations from point clouds whose sampling varies with range and visibility. Bird's eye view (BEV) representations make this problem accessible to efficient two-dimensional convolutional networks. PIXOR established dense, proposal-free oriented detection from BEV inputs~\cite{yang2018pixor}, and PointPillars showed that learned pillar features can support fast detection with a two-dimensional backbone~\cite{lang2019pointpillars}. These approaches motivate a compact architecture that retains useful vertical statistics while allocating most learned processing to the horizontal plane.

Reducing parameter count leaves several design questions unresolved. A richer representation can preserve information discarded by a simple occupancy grid, but its benefit depends on the downstream model. Multi-scale context can strengthen feature processing, yet the resulting detections also depend on the regression objective and their ranking scores. A classification score alone does not explicitly encode localization quality. Conversely, a quality score learned from a footprint does not measure vertical overlap when the detector predicts full three-dimensional boxes. These relationships make component interactions relevant to the design of a small detector.

IRBGHR-PIXOR combines inverted residual blocks, Gaussian heatmap regression, and a modified focal objective for indoor pedestrian perception~\cite{nguyen2024indoor}. The present study addresses outdoor, three-class KITTI detection and develops a different architecture and geometric objective. Its Car, Pedestrian, and Cyclist results therefore answer a different evaluation question from IRBGHR-PIXOR's JRDB pedestrian results. We retain the useful connection between compact convolutions and Gaussian center supervision, while examining how geometry, quality estimation, and contextual refinement work together.

We introduce MobilePixorNeXt with a deterministic fourteen-channel representation, a convolutional hierarchy using $7\times7$ depthwise spatial mixing, scale-gated top-down fusion, and task groups for Car and Pedestrian/Cyclist. Oriented geometric alignment (OGA) combines a multi-axis projection objective with a footprint-scale-normalized corner distance and adaptive task weighting. A separate quality head, referred to as IQA in the implementation, predicts a detached geometric similarity and participates in detection scoring. We investigate these choices through an explicit $2\times2$ OGA--IQA comparison, matched OGA--Q-OGA retraining without a quality head, context and attention comparisons, checkpoint sensitivity, and runtime decomposition.

The study also evaluates a trained three-dimensional extension with shared head convolutions and a vertical regression branch. This extension provides a reference-evaluator result in addition to the local BEV ablation study. A conditional analysis of fixed BEV matches examines why adequate footprints can fail three-dimensional overlap thresholds. The principal contributions are the compact detector implementation, the observed interaction between geometric training and explicit quality prediction, and an evidence-based account of vertical localization limits. All reported runs use one random seed, and the numerical comparisons are confined to the protocols and training budgets under which they were obtained.

\begin{figure*}[t]
\centering
\includegraphics[width=.99\textwidth]{architecture_bev.pdf}
\caption{MobilePixorNeXt architecture used by the P16 BEV ablation. The fourteen-channel input produces a $32\times200\times176$ feature map through three convolutional stages, residual SimAM at stride four, focal context at stride eight, and scale-gated fusion. Car and Pedestrian/Cyclist have separate task groups with independent prediction branches. The trained Final3D extension shares the first head convolution within each group and adds a two-channel vertical branch; it is described in Section~\ref{sec:heads}. Dimensions omit the batch axis.}
\label{fig:architecture}
\end{figure*}
